% Options for packages loaded elsewhere
\PassOptionsToPackage{unicode}{hyperref}
\PassOptionsToPackage{hyphens}{url}
%
\documentclass[
]{article}
\usepackage{amsmath,amssymb}
\usepackage{iftex}
\ifPDFTeX
  \usepackage[T1]{fontenc}
  \usepackage[utf8]{inputenc}
  \usepackage{textcomp} % provide euro and other symbols
\else % if luatex or xetex
  \usepackage{unicode-math} % this also loads fontspec
  \defaultfontfeatures{Scale=MatchLowercase}
  \defaultfontfeatures[\rmfamily]{Ligatures=TeX,Scale=1}
\fi
\usepackage{lmodern}
\ifPDFTeX\else
  % xetex/luatex font selection
\fi
% Use upquote if available, for straight quotes in verbatim environments
\IfFileExists{upquote.sty}{\usepackage{upquote}}{}
\IfFileExists{microtype.sty}{% use microtype if available
  \usepackage[]{microtype}
  \UseMicrotypeSet[protrusion]{basicmath} % disable protrusion for tt fonts
}{}
\makeatletter
\@ifundefined{KOMAClassName}{% if non-KOMA class
  \IfFileExists{parskip.sty}{%
    \usepackage{parskip}
  }{% else
    \setlength{\parindent}{0pt}
    \setlength{\parskip}{6pt plus 2pt minus 1pt}}
}{% if KOMA class
  \KOMAoptions{parskip=half}}
\makeatother
\usepackage{xcolor}
\usepackage[margin=1in]{geometry}
\usepackage{color}
\usepackage{fancyvrb}
\newcommand{\VerbBar}{|}
\newcommand{\VERB}{\Verb[commandchars=\\\{\}]}
\DefineVerbatimEnvironment{Highlighting}{Verbatim}{commandchars=\\\{\}}
% Add ',fontsize=\small' for more characters per line
\usepackage{framed}
\definecolor{shadecolor}{RGB}{248,248,248}
\newenvironment{Shaded}{\begin{snugshade}}{\end{snugshade}}
\newcommand{\AlertTok}[1]{\textcolor[rgb]{0.94,0.16,0.16}{#1}}
\newcommand{\AnnotationTok}[1]{\textcolor[rgb]{0.56,0.35,0.01}{\textbf{\textit{#1}}}}
\newcommand{\AttributeTok}[1]{\textcolor[rgb]{0.13,0.29,0.53}{#1}}
\newcommand{\BaseNTok}[1]{\textcolor[rgb]{0.00,0.00,0.81}{#1}}
\newcommand{\BuiltInTok}[1]{#1}
\newcommand{\CharTok}[1]{\textcolor[rgb]{0.31,0.60,0.02}{#1}}
\newcommand{\CommentTok}[1]{\textcolor[rgb]{0.56,0.35,0.01}{\textit{#1}}}
\newcommand{\CommentVarTok}[1]{\textcolor[rgb]{0.56,0.35,0.01}{\textbf{\textit{#1}}}}
\newcommand{\ConstantTok}[1]{\textcolor[rgb]{0.56,0.35,0.01}{#1}}
\newcommand{\ControlFlowTok}[1]{\textcolor[rgb]{0.13,0.29,0.53}{\textbf{#1}}}
\newcommand{\DataTypeTok}[1]{\textcolor[rgb]{0.13,0.29,0.53}{#1}}
\newcommand{\DecValTok}[1]{\textcolor[rgb]{0.00,0.00,0.81}{#1}}
\newcommand{\DocumentationTok}[1]{\textcolor[rgb]{0.56,0.35,0.01}{\textbf{\textit{#1}}}}
\newcommand{\ErrorTok}[1]{\textcolor[rgb]{0.64,0.00,0.00}{\textbf{#1}}}
\newcommand{\ExtensionTok}[1]{#1}
\newcommand{\FloatTok}[1]{\textcolor[rgb]{0.00,0.00,0.81}{#1}}
\newcommand{\FunctionTok}[1]{\textcolor[rgb]{0.13,0.29,0.53}{\textbf{#1}}}
\newcommand{\ImportTok}[1]{#1}
\newcommand{\InformationTok}[1]{\textcolor[rgb]{0.56,0.35,0.01}{\textbf{\textit{#1}}}}
\newcommand{\KeywordTok}[1]{\textcolor[rgb]{0.13,0.29,0.53}{\textbf{#1}}}
\newcommand{\NormalTok}[1]{#1}
\newcommand{\OperatorTok}[1]{\textcolor[rgb]{0.81,0.36,0.00}{\textbf{#1}}}
\newcommand{\OtherTok}[1]{\textcolor[rgb]{0.56,0.35,0.01}{#1}}
\newcommand{\PreprocessorTok}[1]{\textcolor[rgb]{0.56,0.35,0.01}{\textit{#1}}}
\newcommand{\RegionMarkerTok}[1]{#1}
\newcommand{\SpecialCharTok}[1]{\textcolor[rgb]{0.81,0.36,0.00}{\textbf{#1}}}
\newcommand{\SpecialStringTok}[1]{\textcolor[rgb]{0.31,0.60,0.02}{#1}}
\newcommand{\StringTok}[1]{\textcolor[rgb]{0.31,0.60,0.02}{#1}}
\newcommand{\VariableTok}[1]{\textcolor[rgb]{0.00,0.00,0.00}{#1}}
\newcommand{\VerbatimStringTok}[1]{\textcolor[rgb]{0.31,0.60,0.02}{#1}}
\newcommand{\WarningTok}[1]{\textcolor[rgb]{0.56,0.35,0.01}{\textbf{\textit{#1}}}}
\usepackage{graphicx}
\makeatletter
\newsavebox\pandoc@box
\newcommand*\pandocbounded[1]{% scales image to fit in text height/width
  \sbox\pandoc@box{#1}%
  \Gscale@div\@tempa{\textheight}{\dimexpr\ht\pandoc@box+\dp\pandoc@box\relax}%
  \Gscale@div\@tempb{\linewidth}{\wd\pandoc@box}%
  \ifdim\@tempb\p@<\@tempa\p@\let\@tempa\@tempb\fi% select the smaller of both
  \ifdim\@tempa\p@<\p@\scalebox{\@tempa}{\usebox\pandoc@box}%
  \else\usebox{\pandoc@box}%
  \fi%
}
% Set default figure placement to htbp
\def\fps@figure{htbp}
\makeatother
\setlength{\emergencystretch}{3em} % prevent overfull lines
\providecommand{\tightlist}{%
  \setlength{\itemsep}{0pt}\setlength{\parskip}{0pt}}
\setcounter{secnumdepth}{-\maxdimen} % remove section numbering
\usepackage{booktabs}
\usepackage{caption}
\usepackage{longtable}
\usepackage{colortbl}
\usepackage{array}
\usepackage{anyfontsize}
\usepackage{multirow}
\usepackage{bookmark}
\IfFileExists{xurl.sty}{\usepackage{xurl}}{} % add URL line breaks if available
\urlstyle{same}
\hypersetup{
  pdftitle={Power Calculation for Non-Identifiable Fingermarks Experiment},
  pdfauthor={Maria Cuellar},
  hidelinks,
  pdfcreator={LaTeX via pandoc}}

\title{Power Calculation for Non-Identifiable Fingermarks Experiment}
\author{Maria Cuellar}
\date{2025-05-23}

\begin{document}
\maketitle

\renewcommand{\thetable}{\thesection\arabic{table}} 
\renewcommand{\theHtable}{\thesection\arabic{table}} 
\setcounter{table}{0}

\section{Question of interest}\label{question-of-interest}

I was asked to check whether there are currently sufficient runs for
each group, and if not, how many more are needed for a high-quality
statistical analysis. I approach this by doing an analysis of variance
(ANOVA) test and an ANOVA power calculation.

\section{Data}\label{data}

This is the data from the experiment so far, on 3-4-24.

\begin{table}[t]
\caption*{
{\large Table 1. Original data}
} 
\fontsize{12.0pt}{14.4pt}\selectfont
\begin{tabular*}{\linewidth}{@{\extracolsep{\fill}}lrrrrrrrrrrrrrrrrrr}
\toprule
{\TEXTCOLOR[HTML]{A9A9A9}{TYPE}} & {\TEXTCOLOR[HTML]{A9A9A9}{CALIBER}} & {\TEXTCOLOR[HTML]{A9A9A9}{ROUNDS TESTED}} & {\TEXTCOLOR[HTML]{A9A9A9}{DONORS}} & {\TEXTCOLOR[HTML]{A9A9A9}{4}} & {\TEXTCOLOR[HTML]{A9A9A9}{5}} & {\TEXTCOLOR[HTML]{A9A9A9}{6}} & {\TEXTCOLOR[HTML]{A9A9A9}{7}} & {\TEXTCOLOR[HTML]{A9A9A9}{8}} & {\TEXTCOLOR[HTML]{A9A9A9}{9}} & {\TEXTCOLOR[HTML]{A9A9A9}{10}} & {\TEXTCOLOR[HTML]{A9A9A9}{11}} & {\TEXTCOLOR[HTML]{A9A9A9}{12}} & {\TEXTCOLOR[HTML]{A9A9A9}{13}} & {\TEXTCOLOR[HTML]{A9A9A9}{14}} & {\TEXTCOLOR[HTML]{A9A9A9}{15}} & {\TEXTCOLOR[HTML]{A9A9A9}{16}} & {\TEXTCOLOR[HTML]{A9A9A9}{17}} & {\TEXTCOLOR[HTML]{A9A9A9}{18}} \\ 
\midrule\addlinespace[2.5pt]
{A} & {22} & {198} & {15} & {19} & {11} & {13} & {4} & {1} & {1} & {0} & {1} & {0} & {0} & {0} & {0} & {0} & {0} & {0} \\ 
{A} & {32} & {105} & {12} & {26} & {9} & {12} & {5} & {1} & {3} & {1} & {1} & {0} & {2} & {0} & {0} & {0} & {0} & {0} \\ 
{A} & {38} & {348} & {20} & {38} & {28} & {12} & {11} & {7} & {3} & {0} & {1} & {0} & {0} & {1} & {0} & {0} & {0} & {0} \\ 
{A} & {45} & {178} & {16} & {21} & {16} & {4} & {4} & {3} & {2} & {0} & {0} & {0} & {0} & {0} & {0} & {0} & {0} & {0} \\ 
{R} & {22} & {225} & {23} & {26} & {9} & {9} & {2} & {3} & {1} & {0} & {1} & {0} & {0} & {1} & {0} & {0} & {0} & {0} \\ 
{R} & {38} & {107} & {20} & {30} & {22} & {5} & {6} & {5} & {2} & {5} & {2} & {1} & {0} & {0} & {0} & {0} & {0} & {1} \\ 
{R} & {45} & {102} & {12} & {40} & {17} & {8} & {2} & {9} & {3} & {1} & {1} & {1} & {2} & {0} & {1} & {0} & {0} & {1} \\ 
\bottomrule
\end{tabular*}
\end{table}

There are seven groups defined by the firearm type (Automatic or Rifle)
and the caliber (22, 32, 38, or 45). The groups are A22, A32, A38, A45,
R22, R32, and R45.

We need to rewrite the data so the rows are the individual runs, and the
outcome variable of interest is the number of minutiae found in that
run. Here is an example of the rewritten data.

\begin{table}[t]
\caption*{
{\large Table 2. Sample of rewritten data, first 10 rows}
} 
\fontsize{12.0pt}{14.4pt}\selectfont
\begin{tabular*}{\linewidth}{@{\extracolsep{\fill}}rcr}
\toprule
{\TEXTCOLOR[HTML]{A9A9A9}{RUN}} & {\TEXTCOLOR[HTML]{A9A9A9}{GROUP}} & {\TEXTCOLOR[HTML]{A9A9A9}{MINUTIAE}} \\ 
\midrule\addlinespace[2.5pt]
{1} & {A22} & {3} \\ 
{2} & {A22} & {3} \\ 
{3} & {A22} & {3} \\ 
{4} & {A22} & {3} \\ 
{5} & {A22} & {3} \\ 
{6} & {A22} & {3} \\ 
{7} & {A22} & {3} \\ 
{8} & {A22} & {3} \\ 
{9} & {A22} & {3} \\ 
{10} & {A22} & {3} \\ 
\bottomrule
\end{tabular*}
\end{table}

The table below shows the summary statistics by group. Note that there
are many more observations of A38 than the rest, and the three groups
with the smallest sample sizes are A22, R22, and R45.

\begin{table}[t]
\caption*{
{\large Table 3. Group summary 
 (Mean and variance are of number of minutiae)}
} 
\fontsize{12.0pt}{14.4pt}\selectfont
\begin{tabular*}{\linewidth}{@{\extracolsep{\fill}}crrr}
\toprule
{\TEXTCOLOR[HTML]{A9A9A9}{GROUP}} & {\TEXTCOLOR[HTML]{A9A9A9}{N}} & {\TEXTCOLOR[HTML]{A9A9A9}{MEAN}} & {\TEXTCOLOR[HTML]{A9A9A9}{VARIANCE}} \\ 
\midrule\addlinespace[2.5pt]
{A22} & {50} & {4.30} & {2.13} \\ 
{A32} & {57} & {8.28} & {2.42} \\ 
{A38} & {81} & {11.40} & {2.49} \\ 
{A45} & {23} & {4.04} & {2.68} \\ 
{R22} & {50} & {5.16} & {2.01} \\ 
{R38} & {50} & {9.00} & {1.76} \\ 
{R45} & {54} & {12.17} & {2.07} \\ 
\bottomrule
\end{tabular*}
\end{table}

If we sort the groups, from highest to lowest, by the means, we get:
R45, A38, R38, A32, R22, A22, A45. But we do not know if the differences
between the groups are statistically significant.

\begin{figure}

{\centering \includegraphics[width=0.7\linewidth]{power-calculation-2025-05-23_files/figure-latex/eda-plots-1} 

}

\caption{Comparing minutiae by group.}\label{fig:eda-plots}
\end{figure}

Note that the bottom range is cut out because only the minutiae of 3 and
higher were included. The lower ones were not of interest.

\section{Current power}\label{current-power}

To determine a sufficient sample size for each group, I used an ANOVA
power calculation.

A commonly used rule of thumb for a sufficient power of a study is
between 0.8 and 0.9. If we go with the lower power, we need 20 samples
per group, and if we go with the higher power we need 25. For an ANOVA
test, the same sample sizes for each factor level, i.e., group, should
be equal. Therefore, the groups R45, R22, and A22 do not have sufficient
observations. They each need at least 20.

\section{Data analysis}\label{data-analysis}

\subsection{(NOTE: THIS WILL BE UPDATED WHEN THERE IS MORE
DATA)}\label{note-this-will-be-updated-when-there-is-more-data}

Once we have sufficient data, we can do an ANOVA test to see which of
the groups (firearm type and caliber) has the best acquisition of
minutiae. We should not perform an ANOVA test yet, because there is not
enough data, so the results are meaningless. Nevertheless, as a proof of
concept, the results of an ANOVA test with the current data are the
following.

\begin{Shaded}
\begin{Highlighting}[]
\CommentTok{\# ANOVA test}
\NormalTok{out }\OtherTok{\textless{}{-}} \FunctionTok{aov}\NormalTok{(minutiae }\SpecialCharTok{\textasciitilde{}}\NormalTok{ group, }\AttributeTok{data =}\NormalTok{ dat.long)}
\FunctionTok{summary}\NormalTok{(out)}
\end{Highlighting}
\end{Shaded}

\begin{verbatim}
##              Df Sum Sq Mean Sq F value Pr(>F)    
## group         6   3313   552.2   249.5 <2e-16 ***
## Residuals   358    793     2.2                   
## ---
## Signif. codes:  0 '***' 0.001 '**' 0.01 '*' 0.05 '.' 0.1 ' ' 1
\end{verbatim}

\subsection{Check model assumptions}\label{check-model-assumptions}

The ANOVA test assumes that the data are normally distributed and the
variance across groups are homogeneous. We can check that with some
diagnostic plots.

As the ANOVA test is significant (check this when we have enough
data!!), we can compute Tukey HSD (Tukey Honest Significant Differences,
R function: TukeyHSD()) for performing multiple pairwise-comparison
between the means of groups.

\begin{Shaded}
\begin{Highlighting}[]
\FunctionTok{plot}\NormalTok{(}\FunctionTok{TukeyHSD}\NormalTok{(out))}
\end{Highlighting}
\end{Shaded}

\begin{figure}

{\centering \includegraphics[width=0.7\linewidth]{power-calculation-2025-05-23_files/figure-latex/tukey-test-1} 

}

\caption{Tukey test to determine if there are significant differences between groups, pairwise.}\label{fig:tukey-test}
\end{figure}

\subsubsection{Check the homogeneity of variance
assumption.}\label{check-the-homogeneity-of-variance-assumption.}

\begin{Shaded}
\begin{Highlighting}[]
\FunctionTok{plot}\NormalTok{(out, }\DecValTok{1}\NormalTok{)}
\end{Highlighting}
\end{Shaded}

\begin{figure}

{\centering \includegraphics[width=0.7\linewidth]{power-calculation-2025-05-23_files/figure-latex/aov-diatnostics-1} 

}

\caption{Diagnostic plots for ANOVA test.}\label{fig:aov-diatnostics}
\end{figure}

We see that there are several outliers (points 155, 55, and 54), and
that the residuals are not symmetrical about zero.

Thus, we recommend Levene's test, which is less sensitive to departures
from normal distribution. The function leveneTest() {[}in car package{]}
will be used.

\begin{Shaded}
\begin{Highlighting}[]
\FunctionTok{library}\NormalTok{(car)}
\FunctionTok{leveneTest}\NormalTok{(minutiae }\SpecialCharTok{\textasciitilde{}}\NormalTok{ group, }\AttributeTok{data =}\NormalTok{ dat.long)}
\end{Highlighting}
\end{Shaded}

\begin{verbatim}
## Levene's Test for Homogeneity of Variance (center = median)
##        Df F value Pr(>F)
## group   6  1.0844 0.3712
##       358
\end{verbatim}

\subsubsection{Relaxing the homogeneity of variance
assumption}\label{relaxing-the-homogeneity-of-variance-assumption}

Since the variances by group are different, we perform a Welch one-way
test.

\begin{Shaded}
\begin{Highlighting}[]
\CommentTok{\# out.welch \textless{}{-} oneway.test(minutiae \textasciitilde{} group, data = dat.long) \# not enough data}
\end{Highlighting}
\end{Shaded}

\subsubsection{Check the normality
assumption}\label{check-the-normality-assumption}

\begin{Shaded}
\begin{Highlighting}[]
\FunctionTok{plot}\NormalTok{(out, }\DecValTok{2}\NormalTok{)}
\end{Highlighting}
\end{Shaded}

\begin{figure}
\centering
\pandocbounded{\includegraphics[keepaspectratio]{power-calculation-2025-05-23_files/figure-latex/unnamed-chunk-3-1.pdf}}
\caption{Normality assumption for ANOVA test.}
\end{figure}

\subsubsection{New test if assumptions are not
met}\label{new-test-if-assumptions-are-not-met}

If assumptions are not met, we can use the Kruskal-Wallis rank sum test.

\begin{Shaded}
\begin{Highlighting}[]
\FunctionTok{kruskal.test}\NormalTok{(minutiae }\SpecialCharTok{\textasciitilde{}}\NormalTok{ group, }\AttributeTok{data =}\NormalTok{ dat.long)}
\end{Highlighting}
\end{Shaded}

\begin{verbatim}
## 
##  Kruskal-Wallis rank sum test
## 
## data:  minutiae by group
## Kruskal-Wallis chi-squared = 287.25, df = 6, p-value < 2.2e-16
\end{verbatim}

We do not have enough evidence to suggest that the groups are
statistically significantly different from each other.

\subsection{Analyzing Automatic vs.~Rifle
results}\label{analyzing-automatic-vs.-rifle-results}

Is the outcome different by firearm type? It's not clear from the
boxplot.

\begin{Shaded}
\begin{Highlighting}[]
\NormalTok{dat.long}\SpecialCharTok{$}\NormalTok{type }\OtherTok{\textless{}{-}} \FunctionTok{substr}\NormalTok{(dat.long}\SpecialCharTok{$}\NormalTok{group, }\DecValTok{1}\NormalTok{, }\DecValTok{1}\NormalTok{)}

\NormalTok{dat.long }\SpecialCharTok{\%\textgreater{}\%} \FunctionTok{ggplot}\NormalTok{(}\FunctionTok{aes}\NormalTok{(}\AttributeTok{x=}\NormalTok{type, }\AttributeTok{y=}\NormalTok{minutiae)) }\SpecialCharTok{+} 
  \FunctionTok{geom\_boxplot}\NormalTok{() }\SpecialCharTok{+} 
  \FunctionTok{ylim}\NormalTok{(}\DecValTok{0}\NormalTok{,}\DecValTok{13}\NormalTok{) }\SpecialCharTok{+}
  \FunctionTok{xlab}\NormalTok{(}\StringTok{"Firearm type"}\NormalTok{) }\SpecialCharTok{+} 
  \FunctionTok{ylab}\NormalTok{(}\StringTok{"Minutiae"}\NormalTok{) }\SpecialCharTok{+} 
  \FunctionTok{ggtitle}\NormalTok{(}\StringTok{"Minutiae by firearm type"}\NormalTok{)}
\end{Highlighting}
\end{Shaded}

\begin{figure}

{\centering \includegraphics[width=0.7\linewidth]{power-calculation-2025-05-23_files/figure-latex/type-eda-1} 

}

\caption{Comparing outcome by firearm type.}\label{fig:type-eda}
\end{figure}

We can run a t-test to check whether the Automatic vs.~Rifle results are
significantly different from each other.

\begin{Shaded}
\begin{Highlighting}[]
\NormalTok{out.t }\OtherTok{\textless{}{-}} \FunctionTok{lm}\NormalTok{(minutiae }\SpecialCharTok{\textasciitilde{}} \FunctionTok{as.factor}\NormalTok{(type), }\AttributeTok{data =}\NormalTok{ dat.long)}
\end{Highlighting}
\end{Shaded}

Diagnostic plots for the t test
(\url{https://www.datanovia.com/en/lessons/t-test-in-r/}).

\begin{Shaded}
\begin{Highlighting}[]
\NormalTok{resids }\OtherTok{\textless{}{-}} \FunctionTok{resid}\NormalTok{(out.t)}
\FunctionTok{par}\NormalTok{(}\AttributeTok{mfrow =} \FunctionTok{c}\NormalTok{(}\DecValTok{2}\NormalTok{,}\DecValTok{2}\NormalTok{))}
\FunctionTok{hist}\NormalTok{(resids)}
\FunctionTok{plot}\NormalTok{(out.t, }\AttributeTok{which =} \DecValTok{2}\NormalTok{)}
\FunctionTok{plot}\NormalTok{(out.t, }\AttributeTok{which =} \DecValTok{1}\NormalTok{)}
\FunctionTok{plot}\NormalTok{(out.t, }\AttributeTok{which =} \DecValTok{5}\NormalTok{)}
\end{Highlighting}
\end{Shaded}

\begin{figure}

{\centering \includegraphics[width=1\linewidth]{power-calculation-2025-05-23_files/figure-latex/ttest-diag-1} 

}

\caption{Diagnostic plots for t-test.}\label{fig:ttest-diag}
\end{figure}

We can use a hierarchical model for this to study joint effects of
firearm type and caliber.

\section{Answers for David}\label{answers-for-david}

\subsection{Do we need to collect more samples of the current
calibers?}\label{do-we-need-to-collect-more-samples-of-the-current-calibers}

Yes. We need at least 20 samples for each group, so we need more for
A22, R22, and R45. For the ANOVA test, ideally we would have all the
groups have the same sample size, but this is not necessary.

More specifically, if you have equal variances in the groups and unequal
sample sizes, tnis is not a problem. If you have unequal variances and
equal sample sizes, no problem. But if you have unequal variances and
unequal sample sizes, this is a problem. We might have to run an
alternative test. But this is all ok. The most important thing is
getting at least 20 samples per group.

\subsection{Should we collect R32
data?}\label{should-we-collect-r32-data}

There is no need. We can do the analysis be leaving a zero in the cell
for R32.

In the analysis,one should ignore the structural zeroes in any test of
independence or association. The way to do this is to estimate the
appropriate model on the full contingency table (including the
structural zeros) and then subtract the sum of the chi-square values
associated with the zero cells from the total chi-square test. This
generates a reduced chi-square test of independence only for the reduced
contingency table.

\subsection{Questions for David
(2024-05-03)}\label{questions-for-david-2024-05-03}

\begin{enumerate}
\def\labelenumi{\arabic{enumi}.}
\item
  I assume that the number of minutiae are per Non-Identifiable
  Fingermarks (NIFM) that was found.
\item
  Is there a column for how many NIFMs were found?
\item
  Is it better to find more prints with fewer minutiae each, or fewer
  prints with more minutiae each? Or do we only care about the total
  number of minutiae, regardless of the number of prints?
\item
  What is the outcome of interest?
\end{enumerate}

\subsection{Answers (2024-05-07)}\label{answers-2024-05-07}

\begin{enumerate}
\def\labelenumi{\arabic{enumi}.}
\item
  Yes.
\item
  Yes, in group summary.
\item
  Fewer prints with more minutiae each is better. No, we care about the
  number of prints. There is an interesting analysis that we should do:
  we should give more weight to very informative NIFMs vs.~uninformative
  NIFMs.
\item
  Outcome of interest: How often (among the number of tested rounds) do
  we get a round with a NIFM above a given threshold?
\end{enumerate}

\end{document}
